% BHU PYQ -- Manually transcribed & error-corrected
% Paper: BCF-214 : Corporate Auditing
% Exam: B.Com. (Hons.) F.M.M. Semester III Examination, 2023-24

\documentclass[11pt,a4paper]{article}
\usepackage[a4paper,top=2.5cm,bottom=2.5cm,left=2.8cm,right=2.8cm,headheight=14pt]{geometry}
\usepackage{amsmath,amssymb}
\usepackage[T1]{fontenc}
\usepackage[utf8]{inputenc}
\usepackage{lmodern,microtype,fancyhdr,enumitem,hyperref,lastpage,parskip,booktabs}

\hypersetup{
    colorlinks=true,
    urlcolor=black,
    linkcolor=black,
    pdftitle={BCF-214: Corporate Auditing -- BHU B.Com. (Hons.) FMM Sem III 2023-24}
}

\pagestyle{fancy}
\fancyhf{}
\renewcommand{\headrulewidth}{0.4pt}
\renewcommand{\footrulewidth}{0.4pt}
\fancyhead[L]{\small   \textit{Banaras Hindu University}}
\fancyhead[R]{\small   \textit{BCF-214: Corporate Auditing}}
\fancyfoot[C]{\small Page  \thepage\ of \pageref{LastPage}}


\newlist{parts}{enumerate}{2}
\setlist[parts,1]{label=(\alph*),leftmargin=*,itemsep=5pt,topsep=3pt}
\setlist[parts,2]{label=(\roman*),leftmargin=*,itemsep=3pt,topsep=2pt}

\newcommand{\pts}[1]{\hfill   \textit{[#1]}}


\begin{document}


\begin{center}
  {\large\bfseries BANARAS HINDU UNIVERSITY}\\[2pt]
  {\normalsize B.Com. (Hons.) F.M.M. --- Semester III Examination, 2023-24}\\[6pt]
  \rule{0.8\linewidth}{0.5pt}\\[6pt]
  {\large\bfseries Paper No. BCF-214: Corporate Auditing}\\[4pt]
  {\normalsize Time: 3 Hours\hfill Maximum Marks: 70}
\end{center}

\rule{\linewidth}{0.4pt}

\smallskip



\noindent   \textit{Instructions: Answer all questions. All questions carry equal marks (14 marks each).}

\bigskip




\noindent   \textbf{Question 1.}
What do you mean by auditing? Discuss the basic principles of auditing.\pts{4+10=14}

\centerline{   \textbf{OR}}

What is a continuous audit? Discuss the limitations of continuous audit. Distinguish between continuous audit and periodical audit.\pts{4+5+5=14}

\medskip\rule{\linewidth}{0.2pt}




\noindent   \textbf{Question 2.}
What is an audit notebook? What are the contents of an audit notebook? How does it differ from working papers?\pts{4+5+5=14}

\centerline{   \textbf{OR}}

Explain the steps that an auditor should take before commencement of a new audit.\pts{14}

\medskip\rule{\linewidth}{0.2pt}




\noindent   \textbf{Question 3.}
Define Vouching. What are the objectives of Vouching? Distinguish clearly between the terms `Vouching', `Verification' and `Valuation'.\pts{3+4+7=14}

\centerline{   \textbf{OR}}

How would you verify the following (any two):\pts{7+7=14}

\begin{parts}
  \item Plant and Machinery
  \item Patents
  \item Trade debtors
\end{parts}

\medskip\rule{\linewidth}{0.2pt}




\noindent   \textbf{Question 4.}
What is the objective of `verification of assets'? How far is an auditor responsible as regards their valuation?\pts{7+7=14}

\centerline{   \textbf{OR}}

Describe the provisions of the Companies Act, 2013 with regards to qualifications and disqualifications of the auditor.\pts{14}

\medskip\rule{\linewidth}{0.2pt}




\noindent   \textbf{Question 5.}
Can a properly appointed auditor be removed before the expiry of his term? If so, explain the procedure of removal.\pts{4+10=14}

\centerline{   \textbf{OR}}

Write short notes on any two of the following:\pts{7+7=14}

\begin{parts}
  \item Errors of omission
  \item Government audit
  \item Liabilities of a company auditor
\end{parts}

\vfill
\rule{\linewidth}{0.4pt}

\begin{center}\small
\href{https://bhu-exam.vercel.app/}{Click here to visit our website}\\[4pt]
\href{https://me-aryan.vercel.app/}{Click here to visit our contributor}\\[4pt]
\href{https://quashan-me.vercel.app/}{Click here to visit our contributor}
\end{center}

\end{document}
